\documentclass[12pt]{article}
\usepackage{graphicx} % Required for inserting images
\usepackage{amsmath}
\usepackage{blindtext}

\title{BdOI Monthly Contest  Round1 Solutions: Problem
Blackout Math}
\date{} %
\setlength{\parindent}{0pt}


\begin{document}

\maketitle


\vspace{0.3cm}

\textbf{\Large Problem Information}

\vspace{0.3cm}

{\large Problem Author: Nabeul Islam}

\vspace{0.3cm}

{\large Problem Preparation: Nabeul Islam}

\vspace{0.3cm}

{\large Editorial: Nabeul Islam}

\vspace{0.3cm}



\vspace{0.5cm}

\textbf{\Large Subtasks hinting key observations}

\vspace{0.3cm}

\textbf{\large Preliminary Observation:}

The sequence is defined as
$$F_1 = A$$$$F_2 = B$$$$F_i = F_{i-1} \oplus F_{i-2} \quad \text{for all } i \ge 3$$

Since $N$ can be as large as $10^{18}$, we clearly cannot simulate the recurrence term by term for the full solution within the time constraints. We need to find a efficient solution  to get a full score. However a basic simution will award you the partial points.

\vspace{0.3cm}

\textbf{\large Subtask 1: $N \le 3$}
\vspace{0.3cm}

This is a very easy subtask. For $N \le2$ we just need need print the first and second number in the sequence depending on the input of $N$. For $N=3$ we just need to print $A \oplus B$.


\vspace{0.3cm}

\textbf{\large Subtask 2: $A = 1$ and $B = 1$ and $N \le 10$}

\vspace{0.3cm}

In this subtask , we can simulate $N^{th}$ term in the problem by writing a recursive function. 

\vspace{0.3cm}

\textbf{\large Subtask 3: $N \le 20$}
\vspace{0.3cm}

same as subtask 2.

\vspace{0.3cm}

\textbf{\large Subtask 4:  $N \le 1000$}

\vspace{0.3cm}

In this subtask , we can simulate $N^{th}$ term in the problem by writing a  recursive function. However, it is not efficient always. It can cause stack overflow and give you error. A simple approach to use a loop to simulate the sequence.


\vspace{0.3cm}

\textbf{\large Subtask 5: $N \le 10^5$}
\vspace{0.3cm}

same as subtask 4.

\vspace{0.3cm}


\vspace{0.3cm}

\textbf{\large Subtask 6(Full solution}
\vspace{0.3cm}

Here $N$ is far too large to simulate directly. If we do some dry runs we will find a pattern. The key property to notice here the functionality and property of the XOR($\oplus$) operator. Two of its important properties we need to notice are $x \oplus x = 0$ and $x \oplus 0 = x$.
Let's use this to compute a few terms  and see what happens:

$$F_3 = F_2 \oplus F_1 = A \oplus B$$
$$F_4 = F_3 \oplus F_2 = (A \oplus B) \oplus B = A \oplus (B \oplus B) = A$$
$$F_5 = F_4 \oplus F_3 = A \oplus (A \oplus B) = (A \oplus A) \oplus B = B$$
$$F_6 = F_5 \oplus F_4 = B \oplus A = A \oplus B$$
$$F_7 = F_6 \oplus F_5 = (A \oplus B) \oplus B = A \oplus (B \oplus B) = A$$

Turns out the There is a repeated sequence throughout the terms, we can take this advantange to Optimize our solution.

This pattern repeats in every 3 sequence! 

The sequence is periodic with period 3, cycling through $A, B\documentclass[12pt]{article}
\usepackage{graphicx} % Required for inserting images
\usepackage{amsmath}
\usepackage{blindtext}

\title{BdOI Monthly Contest  Round1 Solutions: Problem
Blackout Math}
\date{} %
\setlength{\parindent}{0pt}


\begin{document}

\maketitle


\vspace{0.3cm}

\textbf{\Large Problem Information}

\vspace{0.3cm}

{\large Problem Author: Nabeul Islam}

\vspace{0.3cm}

{\large Problem Preparation: Nabeul Islam}

\vspace{0.3cm}

{\large Editorial: Nabeul Islam}

\vspace{0.3cm}



\vspace{0.5cm}

\textbf{\Large Subtasks hinting key observations}

\vspace{0.3cm}

\textbf{\large Preliminary Observation:}

The sequence is defined as
$$F_1 = A$$$$F_2 = B$$$$F_i = F_{i-1} \oplus F_{i-2} \quad \text{for all } i \ge 3$$

Since $N$ can be as large as $10^{18}$, we clearly cannot simulate the recurrence term by term for the full solution within the time constraints. We need to find a efficient solution  to get a full score. However a basic simution will award you the partial points.

\vspace{0.3cm}

\textbf{\large Subtask 1: $N \le 3$}
\vspace{0.3cm}

This is a very easy subtask. For $N \le2$ we just need need print the first and second number in the sequence depending on the input of $N$. For $N=3$ we just need to print $A \oplus B$.


\vspace{0.3cm}

\textbf{\large Subtask 2: $A = 1$ and $B = 1$ and $N \le 10$}

\vspace{0.3cm}

In this subtask , we can simulate $N^{th}$ term in the problem by writing a recursive function. 

\vspace{0.3cm}

\textbf{\large Subtask 3: $N \le 20$}
\vspace{0.3cm}

same as subtask 2.

\vspace{0.3cm}

\textbf{\large Subtask 4:  $N \le 1000$}

\vspace{0.3cm}

In this subtask , we can simulate $N^{th}$ term in the problem by writing a  recursive function. However, it is not efficient always. It can cause stack overflow and give you error. A simple approach to use a loop to simulate the sequence.


\vspace{0.3cm}

\textbf{\large Subtask 5: $N \le 10^5$}
\vspace{0.3cm}

same as subtask 4.

\vspace{0.3cm}


\vspace{0.3cm}

\textbf{\large Subtask 6(Full solution)}
\vspace{0.3cm}

Here $N$ is far too large to simulate directly. If we do some dry runs we will find a pattern. The key property to notice here the functionality and property of the XOR($\oplus$) operator. Two of its important properties we need to notice are $x \oplus x = 0$ and $x \oplus 0 = x$.
Let's use this to compute a few terms  and see what happens:

$$F_3 = F_2 \oplus F_1 = A \oplus B$$
$$F_4 = F_3 \oplus F_2 = (A \oplus B) \oplus B = A \oplus (B \oplus B) = A$$
$$F_5 = F_4 \oplus F_3 = A \oplus (A \oplus B) = (A \oplus A) \oplus B = B$$
$$F_6 = F_5 \oplus F_4 = B \oplus A = A \oplus B$$
$$F_7 = F_6 \oplus F_5 = (A \oplus B) \oplus B = A \oplus (B \oplus B) = A$$

Turns out the There is a repeated sequence throughout the terms, we can take this advantange to Optimize our solution.

This pattern repeats in every 3 sequence! 

The sequence is periodic with period 3, cycling through $A, B, A \oplus B,
A, B, A \oplus B, \dots$ forever. 

With this simple oveservation our problem simplifies to the idea below:

 
$$F_N = \begin{cases}
 A & \text{if } N \bmod 3 = 1 \\ 
 B & \text{if } N \bmod 3 = 2 \\ 
 A \oplus B & \text{if } N \bmod 3 = 0 
\end{cases}$$

 This reduces the time complexity and gives a solution that score full points. 
 

\vspace{0.3cm}


\end{document}, A \oplus B,
A, B, A \oplus B, \dots$ forever. 

With this simple oveservation our problem simplifies to the idea below:

 
$$F_N = \begin{cases}
 A & \text{if } N \bmod 3 = 1 \\ 
 B & \text{if } N \bmod 3 = 2 \\ 
 A \oplus B & \text{if } N \bmod 3 = 0 
\end{cases}$$

 This reduces the time complexity and gives a solution that score full points. 
 

\vspace{0.3cm}


\end{document}